\chapter{Expliquer : préparer le terrain d'AIDA}
\label{ch:xai}

\questionchapitre{Une prévision de niveau de nappe n'est mobilisable que si l'on peut dire sur
quoi elle repose et ce qu'elle deviendrait dans d'autres conditions. Comment produire de telles
explications sans qu'elles soient elles-mêmes trompeuses ?}

\section{Pourquoi l'explicabilité commande le reste}

Ce chapitre expose la finalité du travail. L'infrastructure des chapitres précédents a été
construite pour qu'il devienne abordable.

L'exigence d'explication ne relève pas d'une préférence méthodologique. Les niveaux de nappe
conditionnent des décisions administratives contraignantes, restrictions d'usage, autorisations
de prélèvement, dimensionnement d'ouvrages d'alimentation. Une décision publique doit pouvoir
être motivée : un modèle dont on ne sait pas dire pourquoi il annonce une baisse ne sera pas, et
ne devrait pas être, invoqué à l'appui d'un arrêté.

S'y ajoute une condition d'adoption. L'interlocuteur d'un modèle de nappe est un hydrogéologue
qui possède une compréhension physique du système. Il n'accordera sa confiance qu'à un modèle
dont les explications sont compatibles avec ce qu'il sait. Réciproquement, une explication qui
contredit sa compréhension est une information précieuse : soit le modèle a capté quelque chose,
soit il se trompe, et dans les deux cas il faut le savoir.

Cette exigence est le c\oe ur du projet ANR \textsc{Aida}~\cite{anr_aida}, qui vise des
explications \emph{actionnables}. Trois questions organisent le dispositif construit ici : sur
quoi le modèle s'appuie-t-il, que se passerait-il dans d'autres conditions, et cette chronique
est-elle seulement explicable par le climat.

\section{Sur quoi le modèle s'appuie-t-il : les méthodes d'attribution}

Le premier volet mobilise les familles usuelles d'attribution, chacune répondant à une facette
différente de la question.

\begin{center}
\begin{tabularx}{\textwidth}{@{}lX@{}}
\toprule
\textbf{Famille} & \textbf{Ce qu'elle apporte} \\
\midrule
Importance par perturbation & Mesure la dégradation de la prévision quand une entrée est
brouillée, sans hypothèse sur le modèle \\
Valeurs de Shapley~\cite{lundberg2017shap} & Répartition de la contribution entre entrées sur
des bases axiomatiques, avec une variante attribuant aussi dans le
temps~\cite{bento2021timeshap} \\
Attributions par gradients~\cite{sundararajan2017ig,captum} & Saillance temporelle et gradients
intégrés : quels pas de temps de la fenêtre d'entrée pèsent sur la sortie \\
Poids d'attention & Pour les architectures qui en possèdent, lecture directe de ce que le modèle
regarde \\
Décomposition du signal & Sépare tendance, saisonnalité et résidu, et rapporte la prévision à
ces composantes \\
\bottomrule
\end{tabularx}
\end{center}

Disposer de plusieurs familles dans un même outil ne vise pas l'exhaustivité mais la
\textbf{confrontation}. Des méthodes reposant sur des principes différents et qui convergent
donnent une explication crédible ; des méthodes qui divergent signalent une explication fragile,
information utile qu'une méthode unique aurait masquée.

\section{Que se passerait-il si : les contrefactuels sous contrainte physique}
\label{sec:physcf}

C'est l'apport le plus original de ces travaux, et il répond à un problème formulé avec les
hydrogéologues avant d'être traité.

\subsection{Le problème, tel qu'il a été posé}

Une explication contrefactuelle~\cite{wachter2018counterfactual} répond à la question suivante :
quelle modification minimale des entrées changerait la prévision de telle manière ? Appliquée à
une nappe, quel régime climatique aurait évité l'étiage annoncé ?

Les approches existantes pour les séries temporelles~\cite{ates2021comte,wang2023forecastcf}
optimisent la perturbation des entrées sous contrainte de proximité. Sur des données
climatiques, elles produisent des réponses correctes au sens de l'optimisation et physiquement
absurdes : augmenter la température sans modifier l'évapotranspiration, faire varier la pluie
d'un jour à l'autre de façon incohérente, ou construire un régime saisonnier qui n'existe nulle
part.

La réunion de cadrage de février 2026 énonce ce constat dans les termes du métier. Elle relève
que les métriques standard, validité, proximité et compacité, ne captent pas la physique, que le
risque est de « générer des scénarios mathématiquement valides mais physiquement absurdes », et
elle propose deux critères spécifiques : la plausibilité du bilan de masse et la cohérence
temporelle avec les échelles de réponse des aquifères~\cite{reunion_aida}. La méthode décrite
ci-dessous met en \oe uvre ces deux critères.

\subsection{La solution retenue}

L'approche développée renonce à perturber librement les entrées. Elle optimise à la place un
petit nombre de paramètres interprétables, dont chacun a un sens physique et une plage de
validité.

\begin{center}
\begin{tabularx}{\textwidth}{@{}lllX@{}}
\toprule
\textbf{Paramètre} & \textbf{Plage} & \textbf{Neutre} & \textbf{Interprétation} \\
\midrule
Facteur de pluie, un par saison & \numrange{0,3}{2,0} & \num{1,0} & Régime pluviométrique
saisonnier, de \SI{-70}{\percent} au double \\
Décalage de température & \SIrange{-5}{5}{\celsius} & \num{0} & Réchauffement ou
refroidissement global \\
Résidu d'évapotranspiration & \numrange{-0,03}{0,03} & \num{0} & Écart fractionnaire résiduel \\
Décalage temporel & \SIrange{-30}{30}{\day} & \num{0} & Avance ou retard de la saison \\
\bottomrule
\end{tabularx}
\end{center}

Ces paramètres ne sont pas indépendants. L'évapotranspiration est couplée à la température par
la relation de Clausius-Clapeyron, au taux d'environ \SI{7}{\percent} par degré qui gouverne la
capacité de l'air à contenir de la vapeur d'eau. Élever la température dans le scénario
contrefactuel élève donc mécaniquement la demande évaporative : le scénario reste
thermodynamiquement cohérent \emph{par construction}, et non par pénalisation d'une incohérence
après coup. C'est la traduction du premier critère demandé, la plausibilité du bilan de masse.
La structure saisonnière des facteurs de pluie et le décalage temporel traduisent le second, la
cohérence avec les échelles de réponse naturelles.

L'ensemble reste différentiable, donc optimisable par descente de gradient à travers le modèle
de prévision, quel qu'il soit : une couche d'adaptation permet d'appliquer la méthode aux
architectures du catalogue du chapitre~\ref{ch:apprentissage}. Une variante par optimisation
bayésienne est disponible lorsque le gradient n'est pas exploitable.

\subsection{Une comparaison avec une référence non contrainte}

Contraindre l'espace de recherche améliore l'interprétabilité ; encore faut-il montrer que le
prix payé reste acceptable. Une méthode de référence inspirée des contrefactuels pour séries
temporelles multivariées~\cite{ates2021comte} a donc été implémentée pour la comparaison. Elle
optimise directement la perturbation brute sous contraintes souples, sans couplage physique,
soit un nombre de paramètres libres égal à trois fois la longueur de la fenêtre, contre sept
pour l'approche contrainte.

\begin{remarque}[title={Ce que cette comparaison met à l'épreuve}]
Avec deux à trois ordres de grandeur de paramètres libres en plus, la référence non contrainte
atteindra presque toujours la cible plus facilement. La question n'est donc pas laquelle des
deux optimise mieux, ce qui est réglé d'avance, mais si davantage de degrés de liberté produit
une \emph{meilleure explication}. Le dispositif permet d'en discuter sur pièces plutôt que de le
supposer.
\end{remarque}

\subsection{Une validation par un modèle indépendant}

Un contrefactuel produit par optimisation à travers un modèle appris peut n'être valide que pour
ce modèle : il exploiterait alors une particularité de l'approximation plutôt qu'une propriété
du système. Pour écarter cette possibilité, le scénario est rejoué dans un modèle métier calibré
indépendamment, celui du chapitre~\ref{ch:metier}, et l'on vérifie que l'effet annoncé s'y
retrouve.

Cette validation croisée n'était possible que parce que les deux familles de modèles cohabitent
dans le même outil, sur les mêmes données. C'est un bénéfice direct de l'architecture retenue et
une raison supplémentaire d'avoir commencé par le modèle métier.

\section{Cette chronique est-elle explicable par le climat ?}

La troisième question renverse la perspective. Une chronique piézométrique qu'aucun modèle
climatique ne parvient à reproduire n'est pas nécessairement le signe d'un mauvais modèle : elle
peut porter la trace d'une influence anthropique, un prélèvement, que le forçage climatique ne
contient pas.

L'enjeu est double, et il a été formulé comme tel lors du cadrage. Il s'agit d'abord d'écarter
les stations impropres à l'apprentissage, préoccupation qui rejoint directement la constitution
d'un jeu de données de piézomètres peu influencés. Il s'agit ensuite de produire une information
utile en elle-même, l'annotation automatique des piézomètres selon leur caractère inertiel,
réactif ou influencé figurant parmi les cas d'usage identifiés avec le
BRGM~\cite{reunion_aida}.

Le problème est difficile car non supervisé : on ne dispose pas d'un historique étiqueté des
prélèvements. La démarche retenue croise trois regards indépendants, l'analyse des résidus d'un
modèle métier calibré, une détection de ruptures accompagnée d'une analyse d'écart par
apprentissage, et une lecture par représentation apprise, puis fusionne leurs verdicts. La
couche d'explicabilité y joue un rôle inhabituel : elle ne sert pas à justifier une prédiction
mais à détecter une \emph{dérive dans les attributions}, un changement de ce sur quoi le modèle
s'appuie, qui signale que le régime du système a changé. Le résultat est confronté, lorsqu'elle
existe, à l'information publique sur les ouvrages de prélèvement déclarés.

\begin{acquisouvert}
\textbf{Acquis.} Un dispositif d'explicabilité complet et confrontable, articulant attribution,
contrefactuels et détection d'influence. La contribution originale porte sur les contrefactuels
sous contrainte physique : en réduisant l'espace de perturbation à sept paramètres
interprétables couplés par une relation thermodynamique, la méthode ne produit que des scénarios
formulables dans le langage du domaine, un hiver plus sec de \SI{20}{\percent} et plus chaud
d'un degré, au lieu de perturbations sans signification physique. Elle met en \oe uvre les deux
critères demandés par les hydrogéologues, sa validité est éprouvée par un modèle indépendant et
sa qualité mise en regard d'une référence non contrainte.

\textbf{Ouvert.} L'évaluation quantitative comparée des deux approches, sur un échantillon large
et selon des métriques de qualité d'explication établies, reste à conduire. C'est la suite
immédiate de ces travaux et elle relève du projet \textsc{Aida}.
\end{acquisouvert}
